% ===========================================================================
\section{极大理想与素理想}\label{sec:maximal-prime}

\subsection{单环与域的刻画}

\begin{definition}[平凡理想、单环]
零理想 $\{0\}$ 与整个环 $R$ 称为\keyword{平凡理想}；没有非平凡理想的环
称为\keyword{单环}（类比单群：没有非平凡正规子群的群）。
\end{definition}

\begin{theorem}[含幺交换环是域的刻画]
设 $R$ 为含幺交换环，则
\[
  R\text{ 是域} \iff R\text{ 没有非平凡理想（}R\text{ 是单环）}.
\]
\end{theorem}

\begin{proof}
（$\Rightarrow$）域的理想只要包含非零元 $a$，由 $a$ 可逆，
$1 = a^{-1}a$ 也落入该理想，于是理想等于整个环。

（$\Leftarrow$）任取 $a\in R$，$a\ne 0$。因 $R$ 交换，
$(a) = \{\,ra : r\in R\,\}$。$(a)\ne\{0\}$，而理想只有平凡的两个，故
$(a) = R\ni 1$，即存在 $r$ 使 $ra = 1$，$a$ 可逆。每个非零元可逆，
故 $R$ 是域。
\end{proof}

\subsection{极大理想（Maximal Ideal）}

\begin{definition}[极大理想]
理想 $M\subsetneq R$ 称为\keyword{极大理想}，若任何包含 $M$ 的理想
$I$ 只有 $I = M$ 或 $I = R$ 两种可能。
\end{definition}

\begin{theorem}[极大理想的刻画]
设 $R$ 为含幺交换环，$M\subsetneq R$ 为理想，则
\[
  M\text{ 是极大理想} \iff R/M\text{ 是域}.
\]
\end{theorem}

\begin{proof}
由对应定理，$R/M$ 的理想与 $R$ 中包含 $M$ 的理想一一对应。

（$\Rightarrow$）$M$ 极大，故包含 $M$ 的理想只有 $M$ 与 $R$ 两个，
从而 $R/M$ 没有非平凡理想。又 $R/M$ 交换，由上面的定理，$R/M$ 是域。

（$\Leftarrow$）$R/M$ 是域，由上面的定理它没有非平凡理想，故 $R$ 中
包含 $M$ 的理想只有 $M$ 与 $R$ 两个——这正是 $M$ 极大的定义。
\end{proof}

\subsection{素理想（Prime Ideal）}

\begin{definition}[素理想]
理想 $P\subsetneq R$ 称为\keyword{素理想}，若对任意 $a,b\in R$，
\[
  ab\in P \implies a\in P\ \text{或}\ b\in P.
\]
\end{definition}

\begin{example}
在 $\Z$ 中，素数 $p$ 的倍数构成素理想 $(p)$：若 $ab$ 是 $p$ 的倍数，
则 $a, b$ 中至少一个是 $p$ 的倍数。
\end{example}

\begin{theorem}[素理想的刻画]
设 $R$ 为含幺交换环，$P\subsetneq R$ 为理想，则
\[
  P\text{ 是素理想} \iff R/P\text{ 是整环}.
\]
\end{theorem}

\begin{proof}
商环 $R/P$ 中（零元即 $P$ 本身）：
\[
  (a + P)(b + P) = 0 \iff ab + P = P \iff ab\in P, \qquad
  a + P = 0 \iff a\in P.
\]

（$\Rightarrow$）若 $P$ 是素理想：$(a+P)(b+P) = 0$ 即 $ab\in P$，
由素理想定义 $a\in P$ 或 $b\in P$，即 $a+P = 0$ 或 $b+P = 0$。故 $R/P$
无零因子，是整环。

（$\Leftarrow$）若 $R/P$ 是整环：$ab\in P$ 即 $(a+P)(b+P) = 0$，
由无零因子得 $a+P = 0$ 或 $b+P = 0$，即 $a\in P$ 或 $b\in P$。故 $P$
是素理想。
\end{proof}

\begin{remark}
两个方向的证明其实是同一串等价的正读与倒读：$ab\in P$ 与
“$(a+P)(b+P)=0$”本就是一回事，素理想的定义正是“商环无零因子”的
逐元素写法。
\end{remark}
